% БҮЛЭГ 7 — ЁС ЗҮЙН АСУУДАЛ

\chapter{Ёс зүйн асуудал}
\label{Chapter7}

Энэхүү бүлэгт Монгол хэлний хандлага тодорхойлох NLP системийг хөгжүүлэх,
турших, ашиглах явцад тулгардаг ёс зүйн асуудлуудыг дэлгэрэнгүй авч үзнэ.
Өгөгдлийн эх сурвалжийн зөвшөөрөл, annotation-ы ёс зүй, загварын bias,
нууцлалын хамгаалалт, хос хэрэглээний эрсдэл болон AI хэрэгслийн ил тод
ашиглалт зэрэг зургаан үндсэн хэмжээс дээр үнэлгээ хийж, guideline-ийн Ш-5
шаардлагыг бүрэн хангаж байгааг баталгаажуулна.

%-------------------------------------------------------------------------------
\section{Өгөгдлийн ёс зүй ба зөвшөөрөл}

Судалгаанд ашигласан өгөгдөл нь найман нийтэд нээлттэй эх сурвалжаас бүрдэнэ:
мэдээний сайтуудын (news.mn, gogo.mn, IKON.mn, Medee.mn, Eagle News гэх мэт)
нийтийн коммент хэсэг, E-Mongolia, Zaluusinfo, болон Facebook-ийн нээлттэй
группийн коммент (задаргааг бүлэг~\ref{Chapter4}, Хүснэгт~\ref{tab:data_full_corpus}).
Эдгээр өгөгдөл нь хууль, ёс зүйн хувьд дараах нөхцөлийг хангасан байна:

\begin{itemize}
  \item \textbf{Нийтийн эх сурвалж:} Бүх эх сурвалж нь вэб хуудас дээр
    нийтэд нээлттэй хэвлэгдсэн контент бөгөөд хэрэглэгчид нийтлэхдээ олон
    нийтийн үзэх боломжийг ойлгож зөвшөөрсөн байдалтай. Хувийн мессеж,
    хаалттай групп, нэвтрэх зөвшөөрөл шаардсан хуудасны контентыг цуглуулаагүй.
  \item \textbf{Монгол Улсын хууль:} ``Хувь хүний нууцын тухай хууль''
    (2021 оны шинэчилсэн найруулга) болон ``Цахим гарын үсэг, цахим баримт
    бичгийн тухай хууль''-ийн хүрээнд хувийн шинж чанарыг тодорхойлох
    мэдээллийг (PII --- Personally Identifiable Information) задруулах,
    үнэлэх, загвар сургахад ашиглах нь хориотой. Иймд өгөгдлийг цуглуулмагц
    хэрэглэгчийн нэр, профайлын холбоос, зураг, хаяг, утасны дугаар зэргийг
    шүүн хасч, зөвхөн анонимчлагдсан (anonymized) текст агуулгыг хадгалсан.
  \item \textbf{GDPR-ын зарчмуудаас санаа авсан байдал:} Судалгаа нь Монгол
    улсад хамаарах тул ЕХ-ны GDPR-т нийцэх үүрэггүй боловч академик стандартын
    хувьд GDPR-ын үндсэн зарчмуудаас (зорилгын хязгаарлалт, өгөгдлийн
    минимизаци, хадгалах хугацаа) санаа авч удирдлага болгосон. Өгөгдлийг судалгааны
    зорилгоор хязгаарлан ашиглах бөгөөд арилжааны түгээлт, гуравдагч этгээдэд
    шилжүүлэлт зохиохгүй.
  \item \textbf{Нийтлэл эзэмшигчийн эрх:} news.mn-ээс авсан агуулгын зохиогчийн
    эрхийг хүндэтгэн, уг эх сурвалжийг эш татаж, зөвхөн сургалтын вектор болгон
    ашиглав. Бичвэрийг бүтнээр нь нийтлэхгүй.
\end{itemize}

Эдгээр нөхцөл нь academic ethics review-ийн стандартыг хангаж байгаа бөгөөд
IRB (Institutional Review Board)-ийн хянан шийдвэрлэх шаардлагаас ангид
(public, de-identified data) гэж үнэлэв.

%-------------------------------------------------------------------------------
\section{Аннотаци ба аннотаторуудын хоорондын нийцлийн шинжилгээ}

Өгөгдлийг 4 ангилалд шошголохдоо нэг хүний субъектив дүгнэлт хараат байдлыг
бууруулахын тулд олон хүний annotator-ын багийг ашигласан (бүлэг~\ref{Chapter2}).
Багт Монгол хэлний мэдрэмжтэй, академик түвшинд NLP болон хэл шинжлэлд сургагдсан
$n = 3$ annotator хамрагдав. Ёс зүйн шинжтэй дараах процессыг
баримталсан:

\begin{enumerate}
  \item \textbf{Annotation удирдамж (guideline):} Ангилал тус бүрийн
    тодорхойлолт, жишээ, ялгарах шинжийг багтаасан бичгийн удирдамжийг
    бэлтгэж, бүх annotator-д тараав. Маргаантай хилийн тохиолдлуудыг (edge
    case) тусгайлан баримтжуулсан.
  \item \textbf{Сургалтын үе (training phase):} Эхний 200 жишээг бүх
    annotator хамтад нь шошголж, зөрүүтэй шошгуудыг хэлэлцэж удирдамжийг
    боловсруулсан.
  \item \textbf{Бие даасан шошгололт:} Удирдамжийг тохируулсны дараа гурван
    annotator бүх $10{,}000$ дээжийг бие биенээсээ хамааралгүй, бие даан
    шошголсон.
  \item \textbf{Тохиролцооны метрик:} Гурван annotator бүх өгөгдлийг
    шошголсон тул нийцлийг Fleiss' $\kappa$ коэффициентээр хэмжсэн (хоёр
    annotator-т зориулсан Cohen's $\kappa$ биш). Эцсийн утга:
    $\kappa\!=\!0.72$, энэ нь \textit{substantial} нийцэл (0.61--0.80
    хүрээ)-д хамрагдах бөгөөд датасетын найдвартай байдлыг хангаж байна.
  \item \textbf{Маргаан шийдвэрлэх журам:} Зөрүүтэй шошгуудыг гурван
    annotator-ын хамтарсан хэлэлцүүлгээр (консенсус) шийдвэрлэв. Санал
    бүрэн нийлэхгүй, утга тодорхойгүй жишээг ``хилийн тохиолдол''
    (ambiguous) гэж тэмдэглэн өгөгдлийн багцаас хассан.
\end{enumerate}

Inter-annotator agreement-ыг ил тод тайлагнах нь уг датасетын найдвартай
байдлыг батлах, бусад судлаачдад нөхцөл байдлыг ойлгох боломж олгоход
чухал юм. Мөн ангилал тус бүрийн хувьд тохиролцоог тусад нь тайлагнасан:
Хортой сөрөг ба Бүтээлч шүүмжлэлийн хоорондын тохиролцоо нь Эерэг/Саармаг-ээс бага
байх хандлагатай бөгөөд энэ нь эдгээр хоёр ангилалын семантик ойролцоо байдлыг,
улмаар аннотаторуудын субъектив үнэлгээний зөрүүг илэрхийлж байна.

%-------------------------------------------------------------------------------
\section{Bias ба тэгш байдлын шинжилгээ}

Загварын bias (тэгш бус байдал) нь NLP системийн ёс зүйн гол асуудлуудын нэг
юм. Энэхүү систем дараах bias-ын эх сурвалжуудыг үнэлэв:

\subsection{Ангиллын түвшний bias}

Ангиллын тоон тэнцвэргүй байдал нь цөөнх ангиллын F1-д сөрөг нөлөө үзүүлж
магадгүй. Тухайлбал, Эерэг ангилал нь хамгийн цөөн жишээтэй
(бүлэг~\ref{Chapter4}) тул загвар уг ангиллыг ``сурч дуусахгүй'' болж
болзошгүй. Үүнээс урьдчилан сэргийлэхээр:

\begin{itemize}
  \item SMOTE (суурь загварт) ба ангиллын жин оногдуулсан Focal Loss
    ($\gamma{=}2.0$) + WeightedRandomSampler (эцсийн нэг шатлалт MN-BERT-д) хэрэглэв.
  \item Ангилал тус бүрийн F1-ийг тусад нь тайлагнаж, macro F1-ийг анхдагч
    үнэлгээний метрик болгон авав. Accuracy-г дан ганц хэрэглэх нь олонх
    ангилалд хэт итгэмтгий (over-confident) үр дүн өгнө.
\end{itemize}

\subsection{Демографийн bias}

Монгол хэрэглэгчдийн онлайн хамтлаг нь хот/хөдөө, нас, боловсролын түвшин,
улс төрийн чиг баримжааны хувьд олон талт бүрэлдэхүүнтэй. Facebook-оос
цуглуулсан өгөгдөлд залуу, хотын хэрэглэгчид давамгайлдаг тул загвар
хөдөөгийн, ахмад хэрэглэгчдийн хэлний онцлогт дутуу дасан зохицсон байх
эрсдэлтэй. Үүнийг бууруулах арга:

\begin{itemize}
  \item news.mn-ээс авсан өгөгдөл нь илүү албан ёсны, хэвлэмэл Монгол
    хэлэнд ойртсон тул Facebook өгөгдөлтэй хослуулан хэлний хэв маягийн
    олон талт байдлыг нэмэгдүүлсэн.
  \item Эцсийн системийг зөвхөн судалгааны зорилгоор хэрэглэж, бодит
    үйлдвэрлэлийн агуулгын зохицуулалтад (content moderation) ашиглах
    өмнө өргөн хүрээний fairness audit хийх ёстойг баримтжуулсан.
\end{itemize}

\subsection{Бүтээлч шүүмжлэлийг Хортой сөрөг гэж буруу ангилах эрсдэл (false positive)}

Бүлэг~\ref{Chapter6}-ын алдааны шинжилгээнд харуулсан жишээ 1-ийн дагуу,
бүтээлч шүүмжлэлтэй контентыг хортой гэж буруу ангилах нь хэрэглэгчийн хууль ёсны
үзэл бодлыг дарангуйлах эрсдэлтэй. Энэ нь (1) нээлттэй цахим орон зайд ил тод
хэлэлцүүлэг явуулахыг саатуулж, (2) хэвлэл, үзэл бодлоо илэрхийлэх эрхтэй
зөрчилдөж болно.

Үүнийг арилгахаар:

\begin{itemize}
  \item Stage~2-ын Хортой сөрөг ангиллын precision-ыг онцгойлон хянаж, false positive
    rate-ыг хамгийн бага байлгах зорилго тавьсан.
  \item Эцсийн шийдвэрийг хүнд үлдээхээр ``хүн орсон'' (human-in-the-loop)
    архитектурыг нэмэлт хамгаалалт болгон санал болгов.
\end{itemize}

%-------------------------------------------------------------------------------
\section{Нууцлал ба хэрэглэгчийн өгөгдлийн хамгаалалт}

Хэрэглэгчийн нууцлалыг хамгаалах дараах арга хэмжээг баримталсан:

\begin{enumerate}
  \item \textbf{Анонимчлал:} Өгөгдөл цуглуулсны дараа хэрэглэгчийн нэр,
    профайлын ID, профайлын зураг, холбоос зэрэг бүх PII-г шүүн хасав.
    Зөвхөн коммент текст болон хэвлэсэн огноо (жил, сар хүртэл) үлдсэн.
  \item \textbf{Ямар ч PII хадгалдаггүй:} Төслийн бүх кодчилогдсон өгөгдлийн
    файл нь зөвхөн цэвэршүүлсэн текст ба шошгыг агуулна. Үүнийг кодын
    repository-д commit хийхээс өмнө автоматжсан script-ээр шалгаж
    баталгаажуулсан.
  \item \textbf{Хадгалалтын хязгаарлалт:} Түүхий (raw) өгөгдлийг зөвхөн
    судлаачийн ашигладаг локал орчинд хадгалах бөгөөд олон нийтэд
    түгээхгүй. Судалгаа дуусах үед түүхий өгөгдлийг устгана гэж
    төлөвлөсөн.
  \item \textbf{Загварын векторжуулалтын нууцлал:} MN-BERT-ийн контекст
    вектор нь онолын хувьд сургалтын нарийвчилсан өгөгдлийг ``санаж''
    авч болох эрсдэлтэй (membership inference). Уг эрсдэлийг бууруулахаар
    сургалтын epoch-ийг хэт олон болгохгүй, dropout-ыг идэвхжүүлсэн
    хэвээр үлдээв.
  \item \textbf{Үр дүнг тайлагнах ёс зүй:} Алдааны шинжилгээний жишээнүүдэд
    (бүлэг~\ref{Chapter6}) бодит комментыг шууд оруулахгүй, харин хэв
    маягийг илэрхийлэх зохиомол өгүүлбэрүүдийг ашиглав.
\end{enumerate}

%-------------------------------------------------------------------------------
\section{Хос хэрэглээний эрсдэл (Dual-use risk)}

Агуулгын зохицуулалтын зориулалттай NLP систем нь хоёр талт шинж чанартай:
эерэг талаас хортой контентыг шүүж хэрэглэгчдийг хамгаална, харин сөрөг
талаас хууль ёсны үзэл бодлыг дарангуйлах, засаг төрийн эсрэг шүүмжлэлийг
хиймлээр дарагдуулах зэвсэг болж болзошгүй юм.

Монгол орны контекстэд дараах эрсдэлүүд онцгой анхаарал татна:

\begin{itemize}
  \item \textbf{Улс төрийн шүүмжлэлийг ``Хортой сөрөг''-т буруу ангилах:} Жишээ 1-д
    дурдсан бүтээлч шүүмжлэлийг буруу ангилах эрсдэл нь улс төр, нийгмийн
    асуудалд хамаарах шүүмжлэл дээр онцгой анхаарагдана. Хэрэв системийг
    платформын засвар гэдэг зохицуулалтад автомат байдлаар хэрэглэвэл
    ардчилсан хэлэлцүүлгийг сааруулах эрсдэл үүснэ.
  \item \textbf{Цөөнх бүлэг, тусгай хэлний хэлбэрийн талаар bias:} Хэлний
    бага тархацтай хэв маягийг (slang, аялгуу) Хортой сөрөг гэж буруу ангилах
    эрсдэл цөөнх бүлгийн сайн санааны илэрхийлэлд сөргөөр нөлөөлнө.
  \item \textbf{Автомат шийдвэрийн эсрэг хариуцлагагүй байдал:} Алгоритмын
    шийдвэрийг ``объектив'' гэж хүлээн авах нь хүний шүүгчийн хариуцлагыг
    дарагдуулж, алдаатай шийдвэрт эсрэг арга хэмжээ авах боломжийг
    хаах эрсдэлтэй.
\end{itemize}

Эдгээр эрсдэлээс урьдчилан сэргийлэхийн тулд дараах зөвлөмжийг өгнө:

\begin{itemize}
  \item Систем нь зөвхөн ``зөвлөгч'' (advisory) үүрэг гүйцэтгэх ба эцсийн
    шийдвэрийг хүн гаргах (human-in-the-loop).
  \item Precision-ыг хамгийн өндөр байлгах (false positive-ыг багасгах) гэж
    алдааны функц (loss function)-ийг тохируулна.
  \item Загварын таамаглалын үндэслэлийг тайлах (explainability) --- attention
    жин, SHAP-ийн харагдуулалт хэрэглэгчид харуулах ёстой.
  \item Хугацаа тутмын аудит --- улирал бүр санамсаргүй түүвэр дээр хүний
    шалгалт хийх.
\end{itemize}

%-------------------------------------------------------------------------------
\section{AI хэрэгслийн ашиглалт}
\label{sec:ai-disclosure}

Guideline-ийн зүйл 7-д заасны дагуу дипломын ажилд ашигласан AI хэрэгслийг
ил тодоор мэдүүлж байна. Ашигласан гол хэрэгслүүд: \textbf{Anthropic Claude}
(Sonnet/Opus хувилбарууд) ба \textbf{Google Gemini}.

\begin{itemize}
  \item \textbf{Кодын хөгжүүлэлтэд (Claude, Gemini):} PyTorch, HuggingFace
    Transformers, \texttt{imbalanced-learn}, Gradio зэрэг нээлттэй эх
    сурвалжтай сангуудын албан ёсны баримт бичгийг иш татсан. Claude
    ассистантыг гол төлөв debugging, refactoring, DirectML-д нийцтэй
    алгоритмын хэрэгжүүлэлт (FocalLoss-ыг \texttt{torch.gather} ашиглан
    дахин зохион байгуулах гэх мэт) болон grid search дамжлагыг
    бүтээхэд ашигласан. Gemini-г санааг харьцуулах, сонголт хийхэд
    зөвлөгөө авах зорилгоор нэмэлт хэрэглсэн. Бүх код нь зохиогчийн
    гараар шалгагдаж, туршигдсан бөгөөд эцсийн хувилбарт өөрийн
    хяналтаар багтсан.
  \item \textbf{Бичгийн туслалцаанд (Claude, Gemini):} Дипломын текстийн
    редакцийн үе шатанд (хэл найруулга, бүтцийн зохион байгуулалт,
    LaTeX-ын техникийн тусламж, томьёоны бүтэц) Claude-ыг үндсэн
    зөвлөх хэрэгсэл болгон ашигласан; Gemini-г зарим тодорхой асуултад
    хариулт өгөх, эх сурвалж хайх зорилгоор ашигласан. Судалгааны
    санаа, арга зүйн сонголт, үр дүнгийн шинжилгээ, дүгнэлт болон
    чухал шийдвэрүүд нь бүхэлдээ зохиогчийн оруулсан бодит хувь нэмэр
    юм. AI хэрэгслийн гаргасан санал дүгнэлтийг шууд хуулах биш, шүүн
    үзэж, зохиогчийн өөрийн үгээр дахин боловсруулан оруулав.
  \item \textbf{Өгөгдлийн цуглуулалт ба шошголт:} Ямар ч AI хэрэгсэлд
    шошгын шийдвэр хийлгээгүй; бүх 4 ангиллын шошго нь хүний
    annotator-ын гараар хийгдсэн.
  \item \textbf{Эцсийн загвар:} MN-BERT-ийн нарийн тохируулга, суурь загваруудын
    сургалт, Focal Loss + cosine scheduler бүхий grid search, мөн
    туршилтын үнэлгээ нь бүгд зохиогчийн локал орчинд (AMD RX 9070 GRE
    12GB, \texttt{torch-directml}) явагдсан. AI ассистант нь загварын
    жин, гиперпараметрын эцсийн сонголт эсвэл тайлан бүртгэлд шууд
    нөлөөлөөгүй --- зөвхөн debugging, код оновчлолд зөвлөгч үүрэгтэй
    оролцсон.
\end{itemize}

Энэхүү ил тод мэдүүлэлт нь академик үнэнч байдлыг (academic integrity)
хадгалахад чухал бөгөөд guideline-ийн зүйл 7-ын шаардлагыг бүрэн хангаж байна.

%-------------------------------------------------------------------------------
\section{Бүлгийн дүгнэлт}

Энэ бүлэгт өгөгдлийн эх сурвалжийн ёс зүй, annotation-ы тогтолцоо, загварын
bias ба тэгш байдал, нууцлалын хамгаалалт, хос хэрэглээний эрсдэл болон AI
хэрэгслийн ил тод ашиглалт гэсэн зургаан хэмжээст ёс зүйн шинжилгээг явуулсан.
Уг шинжилгээ нь guideline-ийн Ш-5 шаардлагыг бүрэн хангахаас гадна, системийг
цаашид ашиглахад анхаарах ёстой эрсдэлийн жагсаалтыг баримтжуулсан. Гол дүгнэлт
нь: (1)~өгөгдлийг public, de-identified хэлбэрт шилжүүлсэн тул хууль, ёс зүйн
зөвшөөрөл баталгаатай; (2)~аннотаторуудын хоорондын нийцэл ($\kappa\!=\!0.72$)-ийг
тайлагнасан тул датасетын найдвартай байдал ил тод; (3)~Бүтээлч шүүмжлэлийг
Хортой сөрөг гэж буруу ангилах эрсдэл нь хүний хяналтгүй автомат хэрэглэлтэд
зөвшөөрөгдөхгүй --- human-in-the-loop хэрэгцээтэй; (4)~ашигласан AI хэрэгслийн
дэлгэрэнгүй мэдүүлэлтийг хэсэг~\ref{sec:ai-disclosure}-д тодорхойлсон.
